% The actual content of the blueprint, used by both the web and print versions.
% Do not include \begin{document}.
%
% Stage 2 blueprint for rung 4 of ADR 0004: the equality case of the
% Frankl--Kupavskii--Swanepoel threshold. It runs the branch lemmas of rung 3's blueprint
% (projects/FKSEdgeThreshold, FKS Theorem 3) once at the budget C(d+2, 2) instead of
% g(d), with the new geometry that this requires.
%
% Nodes inherited from rung 3 are stated in condensed form, with their \lean{} names. Every
% node below is grounded either in Frankl, Kupavskii and Swanepoel, "Embedding graphs in
% Euclidean space", JCTA 171 (2020) 105146, arXiv 1802.03092 (line numbers are to its LaTeX
% source, main.tex), or in this repository's own statement files. The general results live
% in the shared library GraphDimension, which this project imports through Solution.lean;
% their \lean{} names are library declarations. The statement side lives in this
% repository: the inlined Palomar statement
% (EdgeThresholdUniqueness.Standalone.Mathlib.InlineEdgeThresholdUniqueness), its proof
% sibling, and the advertised target in Solution.lean.
%
% \lean{...} names the declaration that proves the node; `leanblueprint checkdecls` (here:
% `lake exe checkdecls` on the extracted names) verifies it exists in the project's
% environment. Every node below is proved, so each carries \leanok.
%
% The corrections from the review of the argument are recorded where they bite: finding
% A36 (the false general form of Lemma R) at Lemma~\ref{lem:two-apices}, and the review
% points of A37 at Lemma~\ref{lem:simplex-attach}, the equality branches of
% Section~\ref{sec:equality-budget}, and Theorem~\ref{thm:uniqueness}.

\chapter{Uniqueness of $K_{d+2}$ at $\binom{d+2}{2}$ edges for $d \ge 6$}

\section{Notation}

$f(d)$ is the least number of edges of a graph that is not realizable in $\mathbb{R}^d$
(Frankl--Kupavskii--Swanepoel, l.~82). Write $m(d) = \binom{d+2}{2}$ and
$g(2) = 3$, $g(3) = 8$, $g(d) = m(d) - 1$ for $d \ge 4$ (l.~343). $\mathbb{S}^{d-1}$ is the
sphere of radius $1/\sqrt 2$ about the origin of $\mathbb{R}^d$ (l.~57). $S(k)$ is the
statement Frankl--Kupavskii--Swanepoel induct on (l.~344): every graph with at most $g(k)$
edges is realizable in $\mathbb{R}^k$, and every graph with at most $g(k)$ edges that avoids
$K_{k+1}$ and $K_{k+2} - K_3$ is realizable on $\mathbb{S}^{k-1}$. ``Contains'' and
``avoids'' always refer to ordinary subgraphs, not induced ones. Throughout,
$H$ is the $(d-1)$-core of the graph under discussion (Lemma~\ref{lem:core}).

\section{Realizability}

\begin{definition}[Realizability in Euclidean space]
  \label{def:representable}
  \lean{EdgeThresholdUniqueness.Standalone.Mathlib.InlineEdgeThresholdUniqueness.UnitDistanceRepresentable, SimpleGraph.UnitDistEmbeddable, EdgeThresholdUniqueness.Bridge.unitDistanceRepresentable_iff}
  \leanok
  A graph $G$ is \emph{realizable in $\mathbb{R}^n$} when there is an injective
  $f : V(G) \to \mathbb{R}^n$ with $\lVert f(u) - f(v)\rVert = 1$ for every edge $uv$
  (Frankl--Kupavskii--Swanepoel l.~55--57). Non-adjacent vertices are unconstrained: a
  non-edge may also land at distance one. The Palomar statement's inlined definition and
  the library's agree by \texttt{Iff.rfl} (\texttt{unitDistanceRepresentable\_iff}).
\end{definition}

\begin{lemma}[Monotonicity]
  \label{lem:monotone}
  \lean{SimpleGraph.UnitDistEmbeddable.of_le, SimpleGraph.UnitDistEmbeddable.comap, SimpleGraph.UnitDistEmbeddable.of_embedding, SimpleGraph.UnitDistEmbeddable.of_iso}
  \leanok
  \uses{def:representable}
  A subgraph of a graph realizable in $\mathbb{R}^n$ is realizable in $\mathbb{R}^n$, and so
  is an injective homomorphic preimage: the placement transports along graph homomorphisms,
  embeddings, and isomorphisms.
\end{lemma}

\begin{lemma}[$K_{d+2}$ has no realization in $\mathbb{R}^d$; FKS Lemma 14]
  \label{lem:complete-lower}
  \lean{EuclideanGeometry.card_le_of_equilateral, SimpleGraph.hasUnitDistDim_completeGraph, EdgeThresholdUniqueness.Bridge.not_unitDistEmbeddable_complete_of_lt, EdgeThresholdUniqueness.Bridge.ncard_edgeSet_completeGraph, EdgeThresholdUniqueness.Bridge.completeGraph_exists_adj}
  \leanok
  \uses{def:representable}
  At most $d + 1$ points of $\mathbb{R}^d$ are pairwise at distance one, so $K_{d+2}$ has no
  realization in $\mathbb{R}^d$. FKS introduction (l.~85) calls it clear; it is what makes
  the equality case non-vacuous, since $K_{d+2}$ has exactly $m(d)$ edges and no isolated
  vertex. This is also what rules out $K_{d+2}$'s supergraphs, which carry more edges.
\end{lemma}

\section{Spherical placements}

Condensed from rung 3; rung 4 reuses each of them inside the equality-budget branches.

\begin{definition}[Realizability on the sphere; FKS Definition 2]
  \label{lem:sphere-orth}
  \lean{SimpleGraph.SphereEmbeddable, SimpleGraph.SphereEmbeddable.iff_orthogonal, SimpleGraph.SphereEmbeddable.toUnitDist, SimpleGraph.SphereEmbeddable.of_le}
  \leanok
  \uses{def:representable}
  A graph is \emph{realizable on $\mathbb{S}^{d-1}$} when it has an injective placement on
  the sphere of radius $1/\sqrt 2$ putting every edge at distance one. For
  $\lvert x\rvert = \lvert y\rvert = 1/\sqrt 2$, one has $\lvert x - y\rvert = 1$ if and
  only if $\langle x, y\rangle = 0$, and a placement on $\mathbb{S}^{d-1}$ is in particular
  a realization in $\mathbb{R}^d$.
\end{definition}

\begin{lemma}[Poles]
  \label{lem:poles}
  \lean{SimpleGraph.SphereEmbeddable.pole, SimpleGraph.SphereEmbeddable.poles}
  \leanok
  \uses{lem:sphere-orth}
  If $G - v - w$ is on $\mathbb{S}^{d-2}$ and $vw$ is not an edge, then $G$ is on
  $\mathbb{S}^{d-1}$: the section sits in a hyperplane, and $v, w$ go to its two poles. The
  one-vertex version, for any vertex, is the same. FKS l.~384--394. The two-vertex form
  needs $v \ne w$ (A37).
\end{lemma}

\begin{lemma}[The $(d-1)$-core]
  \label{lem:core}
  \lean{SimpleGraph.exists_core, SimpleGraph.exists_core_subset, SimpleGraph.exists_core_maximal, SimpleGraph.induce_edgeFinset_card_le, SimpleGraph.ncard_edgeSet_induce_ne, SimpleGraph.ncard_edgeSet_induce_ne_pair}
  \leanok
  Repeatedly delete a vertex of degree at most $d - 2$. The result $H$, the $(d-1)$-core, is
  an induced subgraph of minimum degree at least $d - 1$ (possibly empty) and does not depend
  on the order: it is the maximal vertex set all of whose vertices have more than $d - 2$
  neighbours inside it. The deleted vertices come with an order in which each has at most
  $d - 2$ neighbours among $H$ and the later-deleted vertices --- the order is part of the
  data, since a pruned vertex's original degree can exceed $d - 2$ (A37) --- and deleting
  $v$, or the pair $v, w$, accounts for exactly the edges meeting the deleted vertices.
\end{lemma}

\begin{corollary}[Re-attaching the pruned vertices; FKS Corollary 12]
  \label{cor:degenerate}
  \lean{SimpleGraph.SphereEmbeddable.exists_core, SimpleGraph.SphereEmbeddable.exists_core_subset, SimpleGraph.SphereEmbeddable.of_degenerate}
  \leanok
  \uses{lem:core, lem:sphere-orth}
  If the core $H$ is on $\mathbb{S}^{d-1}$, so is $G$; if $H$ is empty, $G$ is on
  $\mathbb{S}^{d-1}$. The same principle extends a placement of the core to any vertex set
  containing it, which is how the branches below re-attach a deleted vertex to a placed
  $H - v$.
\end{corollary}

\begin{lemma}[Cross-polytope; FKS Lemma 13, repaired as Lemma M]
  \label{lem:cross}
  \lean{SimpleGraph.SphereEmbeddable.of_compl_matching}
  \leanok
  \uses{lem:sphere-orth}
  If $k \le d$ and the complement of $H$, a graph on $d + k$ vertices, has $k$ disjoint
  edges, then $H$ is on $\mathbb{S}^{d-1}$: the $i$-th pair goes to $\pm e_i/\sqrt 2$, the
  other vertices to the unused $e_j$. The written proof treats only $k = d$ (l.~335); the
  uses below are $k \le 3$.
\end{lemma}

\begin{proposition}[FKS Proposition 2]
  \label{prop:maxdeg}
  \lean{SimpleGraph.SphereEmbeddable.of_degree_le, SimpleGraph.exists_partition_maxDegree_le}
  \leanok
  \uses{lem:sphere-orth}
  Let $d \ge 2$. A graph of maximum degree at most $d - 1$ is on $\mathbb{S}^{d-1}$.
  Induction on $d$ through Lov\'asz's 1966 partition lemma (\texttt{exists\_partition\_maxDegree\_le}),
  with bases $d = 2$ and $d = 3$. FKS l.~77, 122--125.
\end{proposition}

\section{Placements in $\mathbb{R}^d$ off the sphere}

\begin{lemma}[Two simplices on a facet; Lemma S]
  \label{lem:two-simplices}
  \lean{SimpleGraph.exists_twoSimplices_placement}
  \leanok
  \uses{def:representable}
  On $d + 2$ vertices with a fixed pair $\{a, b\}$, every graph without the edge $ab$ is
  realizable in $\mathbb{R}^d$: two unit regular $d$-simplices sharing a facet, with apexes
  $a$ and $b$ at distance $\sqrt{2 + 2/d} < 2$. That is the diameter; every other pair is at
  distance one. FKS l.~355--356.
\end{lemma}

\begin{lemma}[Attaching a low-degree tail, $t \le 2$]
  \label{lem:tail}
  \lean{SimpleGraph.UnitDistEmbeddable.extend_tail}
  \leanok
  \uses{lem:two-simplices}
  Let $d \ge 3$ and fix a placement in which all placed pairs are less than $2$ apart, such
  as Lemma~\ref{lem:two-simplices}'s. A graph that adds new vertices and at most $2$ edges,
  each meeting a new vertex, is realizable extending that placement: every new vertex has
  degree at most two, and two placed neighbours are less than $2$ apart, so the unit spheres
  about them meet in a $(d-2)$-sphere. FKS l.~358--359. Lemma~\ref{lem:tail-three} below
  raises the tail to three edges, which is what the equality budget needs.
\end{lemma}

\begin{lemma}[The point off the sphere; Case B's pole]
  \label{lem:basis-apex}
  \lean{EuclideanGeometry.exists_unit_dist_of_orthogonal_basis, EuclideanGeometry.exists_unit_dist_std_orthonormal_basis}
  \leanok
  \uses{lem:sphere-orth}
  Let $e_1, \dots, e_d$ be pairwise orthogonal with $\lvert e_i\rvert = 1/\sqrt 2$. The point
  $x = 2c \sum e_i$ with $4dc^2 - 4c - 1 = 0$ is at distance one from every $e_i$, and
  $\lvert x\rvert^2 = 2dc^2 \ne 1/2$, so $x$ is off $\mathbb{S}^{d-1}$ and differs from every
  point placed on it. FKS l.~368--369.
\end{lemma}

\section{The equality geometry}
\label{sec:equality-geometry}

Three placements that rung 3 left outside its import closure, and that the equality budget
needs. The library proves each; this project imports them through \texttt{Solution.lean}.

\begin{lemma}[Two apexes and a facet vertex; Lemma R, repaired]
  \label{lem:two-apices}
  \lean{EuclideanGeometry.infinite_common_unit_sphere_two_apices}
  \leanok
  \uses{lem:two-simplices}
  Let $d \ge 4$, $\lvert a - b\rvert^2 = 2 + 2/d$, and $\lvert a - c\rvert = \lvert b - c\rvert = 1$.
  The points at distance one from $a$, $b$ and $c$ form an infinite set: the triangle's
  squared circumradius is $1/(4 - \lvert a-b\rvert^2) = d/(2(d-1)) < 1$, and the solution set
  is a sphere of positive radius in a space of dimension $d - 2 \ge 2$ about the
  circumcentre. The other triples of Lemma~\ref{lem:two-simplices} are equilateral
  (squared circumradius $1/3$) and are covered by
  \texttt{infinite\_common\_unit\_sphere\_twoSimplices} in Lemma~\ref{lem:tail-three}.

  The sketch's general form --- pairwise distances at most $D$ with $D^2 < 3$ imply a
  circumradius below one --- is \textbf{false}: $(-\tfrac34, 0), (\tfrac34, 0),
  (0, \tfrac1{10})$ has squared circumradius $(229/80)^2$ with $D^2 = 9/4 < 3$, because an
  obtuse triangle's circumradius is not half its longest side. That is finding A36, which
  replaced the general lemma by the two special radii that the applications actually need.
\end{lemma}

\begin{lemma}[Attaching a three-edge tail; Lemma T]
  \label{lem:tail-three}
  \lean{SimpleGraph.UnitDistEmbeddable.extend_tail_three, SimpleGraph.infinite_common_unit_sphere_twoSimplices}
  \leanok
  \uses{lem:two-apices}
  Let $d \ge 4$ and fix an injective placement of $s$ in which all placed pairs are less
  than $2$ apart and every triple of distinct placed points has an infinite common
  unit-sphere set. A graph that adds new vertices and at most $3$ edges, each meeting a new
  vertex, is realizable extending that placement. On
  Lemma~\ref{lem:two-simplices}'s configuration every triple has an infinite common
  unit-sphere set, by Lemma~\ref{lem:two-apices} for the triples containing both apexes and
  by the equilateral case otherwise. The placement is built one vertex at a time: whenever
  the boundary of placed vertices with unplaced neighbours has size at most three, the
  common unit-sphere set is infinite, so the finitely many placed points can be avoided.
\end{lemma}

\begin{lemma}[Attaching to a simplex; Lemma P]
  \label{lem:simplex-attach}
  \lean{EuclideanGeometry.infinite_sphere_inter_of_regular_simplex, SimpleGraph.UnitDistEmbeddable.extend_of_clique_neighbors}
  \leanok
  Let $K$ be a unit regular $d$-simplex in $\mathbb{R}^d$. New vertices, each joined to a
  set of at most $d - 1$ vertices of $K$ (possibly empty; the sets may overlap), attach
  injectively to a placement of $K$: each solution set is a sphere of positive dimension
  about the sub-simplex's circumcentre (A37), so the finitely many placed points are
  avoided. Neither disjoint neighbourhoods nor the bound $s \le d - 2$ is needed: the
  packaged form \texttt{extend\_of\_clique\_neighbors} asks only that each outside vertex's
  neighbourhood be a clique of size at most $d - 1$, and allows outside vertices adjacent to
  each other. The core \texttt{infinite\_sphere\_inter\_of\_regular\_simplex} covers $k + 1
  \le d$ points of the simplex.
\end{lemma}

\section{The induction, with the budget as a parameter}
\label{sec:induction}

Condensed from rung 3, where the step is proved. Throughout, $d \ge 3$, $G$ is a finite
graph with at most $b$ edges, and $H$ is its $(d-1)$-core. Every branch of the induction
step is a lemma about a graph with at most $b$ edges and a hypothesis on $b$, not about
$g(d)$: rung 3 applies them at $b = g(d)$, rung 4 re-applies the same lemmas at
$b = m(d)$ on top of $S(d-1)$. Nothing in the step is proved twice.

\begin{lemma}[$S(2)$]
  \label{lem:S2}
  \lean{SimpleGraph.fksStatement_two}
  \leanok
  \uses{lem:sphere-orth}
  At most three edges: realizable in $\mathbb{R}^2$. If it also avoids $K_3$ and
  $K_{1,3} = K_4 - K_3$, it is on $\mathbb{S}^1$. FKS: easy to verify.
\end{lemma}

\begin{lemma}[Branch A: the core contains $K_{d+2} - K_3$]
  \label{lem:branch-A}
  \lean{SimpleGraph.UnitDistEmbeddable.of_core_completeMinusTriangle}
  \leanok
  \uses{lem:core, lem:two-simplices, lem:tail}
  Assume $b \le m(d)$ and $b < m(d) + d - 4$. Then $H$ has exactly the $d + 2$ vertices of
  that subgraph, since one more core vertex brings at least $d - 1$ further edges. Then
  either $H = K_{d+2}$, or $H$ misses $s \in \{1, 2, 3\}$ triangle edges, $H$ is placed by
  Lemma~\ref{lem:two-simplices}, and the tail $t$ satisfies $t + m(d) \le b + s$ (stated
  additively; natural subtraction is wrong here, A37). At $b = g(d)$, $t \le s - 1 \le 2$,
  and Lemma~\ref{lem:tail} attaches it. FKS l.~352--359. At $b = m(d)$ the tail reaches
  $3$, which is where Lemma~\ref{lem:tail-three} takes over
  (Lemma~\ref{lem:eq-branch-A}).
\end{lemma}

\begin{lemma}[Branch B: the core contains $K_{d+1}$ and avoids $K_{d+2} - K_3$]
  \label{lem:branch-B}
  \lean{SimpleGraph.UnitDistEmbeddable.of_core_clique, SimpleGraph.UnitDistEmbeddable.of_core_containing_clique}
  \leanok
  \uses{lem:core, cor:degenerate, lem:basis-apex, lem:simplex-attach}
  If $b < m(d) + d - 4$, then $H = K_{d+1}$: two core vertices outside the clique overshoot,
  and one would give $K_{d+2} - K_3$. At most $b - \binom{d+1}{2}$ edges meet the pruned
  vertices. If some clique vertex has no pruned neighbour, the clique-less-$v$ goes on the
  sphere as a scaled orthonormal basis, the pruned vertices are re-attached by
  Corollary~\ref{cor:degenerate}, and $v$ goes to the point of Lemma~\ref{lem:basis-apex}
  (FKS l.~363--369). Otherwise every clique vertex meets the outside, which at the equality
  budget exhausts the tail; that case is Lemma~\ref{lem:eq-core-clique}. The review (A37)
  removed the claim that ``some clique vertex has no outside neighbour'' follows from a
  cross-edge count at $b = m(d)$: it is false there, and unneeded --- the residual case is
  placed directly by Lemma~\ref{lem:simplex-attach}.
\end{lemma}

\begin{lemma}[Branch C, small cores]
  \label{lem:branch-C-small}
  \lean{SimpleGraph.SphereEmbeddable.of_fks_small_core}
  \leanok
  \uses{lem:cross}
  If $H$ avoids $K_{d+1}$ and $K_{d+2} - K_3$ and has at most $d + 2$ vertices, then $H$ is
  on $\mathbb{S}^{d-1}$ by Lemma~\ref{lem:cross} with $k \le 2$. No budget is used.
  FKS l.~373--380.
\end{lemma}

\begin{lemma}[Branch C, a universal vertex]
  \label{lem:branch-C-universal}
  \lean{SimpleGraph.SphereEmbeddable.of_fks_pole}
  \leanok
  \uses{lem:poles}
  Assume $S(d-1)$, $b - (d + 2) \le g(d-1)$, and that $H$ avoids $K_{d+1}$ and
  $K_{d+2} - K_3$, has at least $d + 3$ vertices, and has a vertex $v$ adjacent to all the
  others. Then $H - v$ avoids $K_d$ and $K_{d+1} - K_3$, $S(d-1)$ puts it on
  $\mathbb{S}^{d-2}$, and $v$ goes to a pole. The budget condition holds at $b = g(d)$ for
  every $d \ge 3$, and at $b = m(d)$ with slack zero for $d \ge 5$; at $d = 4$ it fails,
  which is where $K_{3,3,1}$ lives. FKS l.~384--387.
\end{lemma}

\begin{lemma}[Branch C, two poles]
  \label{lem:branch-C-poles}
  \lean{SimpleGraph.SphereEmbeddable.of_fks_poles}
  \leanok
  \uses{lem:poles, prop:maxdeg}
  Assume $S(d-1)$ and $b - (2d - 1) \le g(d-1)$, and that $H$ avoids $K_{d+1}$ and
  $K_{d+2} - K_3$, has at least $d + 3$ vertices, and no vertex of maximum degree is
  universal. If $\Delta(H) \le d - 1$, Proposition~\ref{prop:maxdeg} places $H$. Otherwise,
  for a vertex $v$ of maximum degree and a non-neighbour $w \ne v$ (A37), $H - v - w$ has at
  most $g(d-1)$ edges. If it avoids $K_d$ and $K_{d+1} - K_3$, $S(d-1)$ and
  Lemma~\ref{lem:poles} place $H$. FKS l.~388--391.
\end{lemma}

\begin{lemma}[Case 1 does not occur]
  \label{lem:case1}
  \lean{SimpleGraph.not_forall_clique_of_edgeSet_ncard_lt, SimpleGraph.exists_cliqueFree_pair_of_edgeSet_ncard_lt}
  \leanok
  Assume $b < m(d) + d - 4$, and the hypotheses in force at this point: $H$ avoids
  $K_{d+1}$, has at least $d + 3$ vertices and minimum degree at least $d - 1$,
  $\Delta(H) \ge d$, and each maximum-degree vertex has a non-neighbour $w \ne v$. Then some
  non-neighbour of $v$ leaves a $d$-clique-free deletion: the core cannot be such that every
  deletion $H - v - w$ contains a $d$-clique. Two $d$-cliques meeting in at most $d - 2$
  vertices overshoot by the count $d + m(d) - 4$; meeting in $d - 1$, a further vertex $y$
  overshoots by $\binom{d+1}{2} + 2d - 3$. Pure counting. FKS l.~395--401.
\end{lemma}

\begin{lemma}[Case 2: counting]
  \label{lem:case2-count}
  \lean{SimpleGraph.fks_case_two_degrees}
  \leanok
  If $v$ has maximum degree $\Delta \ge d$ and a non-neighbour $w$ with
  $H - v - w \supseteq K_{d+1} - K_3$, then $\lvert E(H)\rvert \ge \binom{d+1}{2} - 3 +
  \Delta + (d - 1)$, which is $m(d) + d - 5$ at $\Delta = d$. So Case 2 does not occur when
  $b < m(d) + d - 5$: at $b = g(d)$ for $d \ge 5$, and at $b = m(d)$ for $d \ge 6$. At
  equality, $H - v - w = K_{d+1} - K_3$, $\Delta = d$, and $\deg w = d - 1$. FKS
  l.~404--406.
\end{lemma}

\begin{lemma}[Case 2: the equality graphs]
  \label{lem:case2-place}
  \lean{SimpleGraph.SphereEmbeddable.of_case_two, SimpleGraph.SphereEmbeddable.of_fks_case_two}
  \leanok
  \uses{lem:case2-count, lem:cross, lem:sphere-orth}
  In the equality configuration of Lemma~\ref{lem:case2-count}, $H$ is on
  $\mathbb{S}^{d-1}$: either the components of $H[v, w, v_1, v_2, v_3]$ are paths of at most
  three edges, placed on a great circle with the remaining $K_{d-2}$ on the orthogonal
  subsphere; or there are three disjoint non-edges and Lemma~\ref{lem:cross} applies with
  $k = 3$. Rung 3 needs it only at $d = 3$ and $d = 4$; rung 4 at $d \ge 6$ never needs it,
  since $m(d) < m(d) + d - 5$ there. FKS l.~406--412.

  This is the one place the published proof needed a repair. FKS read the equality count
  (l.~404--406) as forcing $H - v - w = K_{d+1} - K_3$, $\deg v = d$, $\deg w = d - 1$, and
  hence $\lvert V(H)\rvert = d + 3$. At $d = 3$ that deduction is false: a vertex joined
  only to $v$ and $w$ has degree $2 = d - 1$ and survives the pruning, and the library
  records a seven-vertex counterexample
  (\texttt{exists\_fks\_case\_two\_extra\_vertex}, \texttt{Extremal/FKS/CaseTwoBoundary.lean},
  Math finding A38; a library leaf outside this project's import closure). The theorem is
  unaffected: at $d = 3$ the equality configuration makes $H - v - w$ a star $K_{1,3}$ plus
  isolated vertices, whose centre $c$ has degree $3 = \Delta$ and is not adjacent to $v$,
  and $H - c - v$ has only two edges, so $S(2)$ and the poles place $H$ --- a
  contradiction. At $d \ge 4$ an extra vertex would need $d - 1 \ge 3$ neighbours among
  $\{v, w\}$, so the deduction holds. The bridge \texttt{SphereEmbeddable.of\_fks\_case\_two}
  (\texttt{Extremal/FKS/CaseTwoBridge.lean}) implements this repair.
\end{lemma}

\begin{lemma}[The step]
  \label{lem:step}
  \lean{SimpleGraph.SphereEmbeddable.of_fks_core, SimpleGraph.SphereEmbeddable.of_core_maximal, SimpleGraph.fksStatement_succ, SimpleGraph.fksBudget_le_pred_add, SimpleGraph.fksBudget_le_pred_add_two_mul, SimpleGraph.fksBudget_lt_choose, SimpleGraph.fksBudget_lt_branch_count}
  \leanok
  \uses{lem:core, cor:degenerate, lem:branch-A, lem:branch-B, lem:branch-C-small,
    lem:branch-C-universal, lem:branch-C-poles, lem:case1, lem:case2-count, lem:case2-place}
  For $d \ge 3$, $S(d-1)$ implies $S(d)$. Take $H$ to be the maximal core. If it is empty,
  Corollary~\ref{cor:degenerate} places $G$; if it is on the sphere,
  Corollary~\ref{cor:degenerate} re-attaches the rest. If $H$ contains $K_{d+2} - K_3$,
  Branch A applies; if it contains $K_{d+1}$, Branch B; otherwise the spherical dispatch runs
  Branch C small, then Proposition~\ref{prop:maxdeg}, then the universal-vertex and
  two-pole branches, with Case 1 excluded by Lemma~\ref{lem:case1} and Case 2 placed through
  the repaired bridge of Lemma~\ref{lem:case2-place}. Arithmetic facts used:
  $g(d) - (d + 2) \le g(d-1)$ and $g(d) - (2d - 1) \le g(d-1)$ for $d \ge 3$, and
  $g(d) < m(d) + d - 4$ for $d \ge 3$. The induction is a two-parameter recursion, over $d$
  and over the core, so this node only dispatches cases.
\end{lemma}

\begin{theorem}[Frankl--Kupavskii--Swanepoel, Theorem 3]
  \label{thm:fks}
  \lean{SimpleGraph.fksStatement, SimpleGraph.unitDistEmbeddable_of_ncard_edgeSet_lt}
  \leanok
  \uses{lem:S2, lem:step, lem:complete-lower}
  For $d \ge 4$, every finite simple graph with fewer than $\binom{d+2}{2}$ edges is
  realizable in $\mathbb{R}^d$; with Lemma~\ref{lem:complete-lower}, $f(d) = \binom{d+2}{2}$.
  FKS l.~91. The library assembles $S(d)$ for every $d \ge 2$ from $S(2)$ and the step;
  rung 4 uses $S(d-1)$ inside the equality-budget spherical branch
  (Lemma~\ref{lem:eq-core-spherical}).
\end{theorem}

\section{The equality budget in the branches}
\label{sec:equality-budget}

The branches of Section~\ref{sec:induction}, re-run at $b = m(d)$ for $d \ge 6$. Every count
in the induction step has excess $d - 4$ over the budget, except Case 2 with $d - 5$ (A23),
so at $b = m(d)$ each branch closes with the new geometry of
Section~\ref{sec:equality-geometry}. Each lemma below states
its budget hypotheses explicitly, in the form the library proves them.

\begin{lemma}[The spherical branch at $b = m(d)$]
  \label{lem:eq-core-spherical}
  \lean{SimpleGraph.SphereEmbeddable.of_fks_core_le_choose}
  \leanok
  \uses{lem:step, lem:branch-C-small, prop:maxdeg, lem:branch-C-universal, lem:branch-C-poles, lem:case1}
  Let $d \ge 6$, and let $G$ have minimum degree at least $d - 1$, at most $m(d)$ edges,
  no $K_{d+1}$ and no $K_{d+2} - K_3$. Then $G$ is on $\mathbb{S}^{d-1}$. The proof is
  Lemma~\ref{lem:step}'s spherical dispatch with $S(d-1)$ in place: small cores by
  Lemma~\ref{lem:branch-C-small}, then Proposition~\ref{prop:maxdeg}, then the two pole
  branches, whose budget conditions read $m(d) - (d+2) \le g(d-1)$, with equality, and
  $m(d) - (2d-1) \le g(d-1)$, with slack; Case 1 is
  excluded because $m(d) < m(d) + d - 4$, and Case 2 cannot arise since the equality count
  overshoots $m(d)$ by $d - 5 > 0$.
\end{lemma}

\begin{lemma}[Branch A at $b = m(d)$: place, or the core is complete]
  \label{lem:eq-branch-A}
  \lean{SimpleGraph.unitDistEmbeddable_or_isNClique_of_core_completeMinusTriangle}
  \leanok
  \uses{lem:branch-A, lem:core, lem:tail-three, lem:two-simplices}
  Let $d \ge 4$ and let $H$ be a core containing $K_{d+2} - K_3$, with
  $b \le m(d)$ (A37 adds this to $b < m(d) + d - 4$) and minimum degree at least $d - 1$.
  Then either $G$ is realizable in $\mathbb{R}^d$, or $H$ is a clique on exactly $d + 2$
  vertices. In the first alternative $H$ misses $s \in \{1, 2, 3\}$ triangle edges and is
  placed by Lemma~\ref{lem:two-simplices}; the tail is counted additively, $t + m(d) \le b + s$
  with $t \le s \le 3$, and Lemma~\ref{lem:tail-three} attaches it --- every triple of
  placed points has an infinite common unit-sphere set, which is what three tail edges need.
\end{lemma}

\begin{lemma}[Branch B at $b = m(d)$: the clique core]
  \label{lem:eq-core-clique}
  \lean{SimpleGraph.UnitDistEmbeddable.of_core_containing_clique_le_choose, SimpleGraph.UnitDistEmbeddable.of_core_clique_le_choose, SimpleGraph.UnitDistEmbeddable.of_core_clique_without_external_neighbor}
  \leanok
  \uses{lem:branch-B, lem:core, cor:degenerate, lem:basis-apex, lem:simplex-attach}
  Three forms, in the order the library layers them. First, a core containing $K_{d+1}$,
  avoiding $K_{d+2} - K_3$, with minimum degree at least $d - 1$ inside it,
  $b \le m(d)$ and $b < m(d) + d - 4$, is exactly $K_{d+1}$: a larger core overshoots the
  count (FKS l.~361). Second, for a maximal core $H = K_{d+1}$ of a graph with at most
  $\binom{d+1}{2} + (d + 1) = m(d)$ edges, $G$ is realizable in $\mathbb{R}^d$. If some
  clique vertex $v$ has no outside neighbour, $K_{d+1} - v$ goes on the sphere,
  Corollary~\ref{cor:degenerate} re-attaches the pruned vertices inside the clique, and $v$
  goes to the off-sphere point of Lemma~\ref{lem:basis-apex}; this arm does not consult the
  tail budget. Otherwise every clique vertex meets the outside: the $d + 1$ cross edges
  exhaust the tail, every outside vertex's neighbours lie in the clique and number at most
  $d - 2$ --- a $(d-1)$-st would make it a core vertex after all --- and
  Lemma~\ref{lem:simplex-attach} places the outside vertices at once. The review (A37)
  removed the earlier claim that the first arm follows from a cross-edge count at this
  budget; both arms are proved, and the branch returns a placement, not a dichotomy.
\end{lemma}

\begin{lemma}[A clique exhausting the budget]
  \label{lem:eq-clique-exhausts}
  \lean{SimpleGraph.nonempty_iso_completeGraph_of_isNClique}
  \leanok
  If $G$ has an $n$-clique on a vertex set $s$, at most $\binom{n}{2}$ edges in total, and
  no isolated vertex, then every edge lies inside $s$: an edge leaving $s$ would exceed the
  budget. So $G = K_n$ on all of $V$, up to isomorphism. This is how
  Branch~A's clique alternative becomes the theorem's conclusion: a $K_{d+2}$ core with
  $m(d)$ edges and no isolated vertex leaves no room for any other edge or vertex.
\end{lemma}

\section{Uniqueness at the complete-graph edge count}

\begin{theorem}[Uniqueness for $d \ge 6$]
  \label{thm:uniqueness}
  \lean{SimpleGraph.nonempty_iso_completeGraph_of_not_unitDistEmbeddable}
  \leanok
  \uses{lem:core, lem:complete-lower, lem:eq-core-spherical, lem:eq-branch-A, lem:eq-core-clique, lem:eq-clique-exhausts}
  Let $d \ge 6$. Every finite simple graph with exactly $\binom{d+2}{2}$ edges, with no
  isolated vertex, and with no unit-distance representation in $\mathbb{R}^d$, is isomorphic
  to $K_{d+2}$. The hypothesis is not vacuous by Lemma~\ref{lem:complete-lower}. The proof
  takes the maximal core $H$. If $H$ contains $K_{d+2} - K_3$, Lemma~\ref{lem:eq-branch-A}
  either places $G$ --- contradicting the hypothesis --- or makes $H$ a $(d+2)$-clique, and
  Lemma~\ref{lem:eq-clique-exhausts} leaves nothing outside it. Otherwise $H$ avoids
  $K_{d+2} - K_3$, and the proof derives the placement that the hypothesis forbids: if $H$
  also avoids $K_{d+1}$, Lemma~\ref{lem:eq-core-spherical} places $G$ on the sphere; if
  $H$ contains $K_{d+1}$, Lemma~\ref{lem:eq-core-clique} places $G$ in $\mathbb{R}^d$.
  Every
  route ends in a placement, so the hypothesis forces the clique alternative. To our
  knowledge this equality case has not previously been stated; the argument is the
  Frankl--Kupavskii--Swanepoel induction run at $b = m(d)$, with the corrections A36 and
  A37 recorded at Lemma~\ref{lem:two-apices}, Lemma~\ref{lem:simplex-attach} and
  Section~\ref{sec:equality-budget}. At $d = 5$ the same induction also closes: Case 2's
  equality graph is impossible by maximum degree, so uniqueness at $d = 5$ needs no new
  coordinates (A37); only $d = 4$ genuinely fails, where $K_{3,3,1}$ escapes.
\end{theorem}

\begin{theorem}[The Palomar statement]
  \label{thm:palomar}
  \lean{EdgeThresholdUniqueness.Standalone.Mathlib.InlineEdgeThresholdUniqueness.UniqueKComplete, EdgeThresholdUniqueness.Standalone.Mathlib.InlineEdgeThresholdUniqueness.UniqueKComplete.proof, EdgeThresholdUniqueness.Standalone.Mathlib.InlineEdgeThresholdUniqueness.UniqueKComplete.witness, EdgeThresholdUniqueness.Palomar.target}
  \leanok
  \uses{thm:uniqueness, def:representable}
  The frozen statement quantifies over $d \ge 6$ and graphs $G$ on \texttt{Fin n}: exactly
  $\binom{d+2}{2}$ edges, no isolated vertex, no unit-distance representation in
  $\mathbb{R}^d$ (Definition~\ref{def:representable}'s inlined predicate), concluding
  $\texttt{Nonempty} (G \simeq_g K_{d+2})$. Its proof sibling specializes
  Theorem~\ref{thm:uniqueness} and reads the inlined predicate as
  \texttt{UnitDistEmbeddable}; \texttt{Solution.lean} restates it for the Comparator as
  \texttt{EdgeThresholdUniqueness.Palomar.target}. The witness records that at $d = 6$ the
  hypotheses hold together: $K_8$ has $28$ edges, no isolated vertex, and no representation
  in $\mathbb{R}^6$.
\end{theorem}
